\begin{figure}[H]
    \centering
    {\chemnameinit{\chemfig{{\color{red}H}-C(=[2]O)-CH_3}}
    \schemestart
        \chemname{\chemfig{{\color{blue}^{-}OOC}-C(=[2]O)-CH_3}}{\small\bfseries Piruvato}
        \arrow(.base east--.base west){-U>[{\color{red}H}][][][0.3][60]}[,2]
        \chemfig{{\color{blue}C} | {\color{blue}O_2}}
        \+
        \chemname{\chemfig{{\color{red}H}-C(=[2]O)-CH_3}}{\small\bfseries Acetaldehido}
    \schemestop
    \chemnameinit{}
    \vspace*{1em}
    \par}
    \rule{\linewidth}{0.2pt}
    \caption[Descarboxilación del piruvato en acetaldehido por la piruvato descarboxilasa]{Reacción catalizada por la piruvato descarboxilasa (EC 4.1.1.1). Esta retira el grupo carboxilo (\ce{COOH}), reacción llamada descarboxilación y de donde proviene el nombre descarboxilasa. Tras esto un átomo de hidrógeno toma el lugar del grupo carboxilo.}
    \label{fig:reaccion_piruvato_descarboxilasa}
\end{figure}